\subsection{有限生成自由模情形下内乘的显式形式}

设 \(\mathbf{M}\) 是一个有限生成的自由 A-模， \((e_i)_{1\leq i\leq n}\) 是 \(\mathbf{M}\) 的一组基，而 \((e_i^*)_{1\leq i\leq n}\) 是 \(\mathbf{M}^*\) 的对偶基。对每个有限序列 \(s = (i_1,\ldots ,i_p)\) （由 \([1,n]\) 中元素组成），设 \(e_s = e_{i_1}\otimes e_{i_2}\otimes \dots \otimes e_{i_p}\) （相应地 \(e_s^* = e_{i_1}^*\otimes \dots \otimes e_{i_p}^*\) ）。我们知道（§ 5，第 5 号，定理 1）， \(e_s\) 构成 A-模 \(\mathsf{T}(\mathbf{M})\) 的一组基，而 \(e_s^*\) 构成 A-模 \(\mathsf{T}(\mathbf{M}^*)\) 的一组基。若 \(s,t\) 是两个有限序列（由 \([1,n]\) 中元素组成），设 \(s.t\) 表示如下得到的序列：如果 \(s = (i_1,\ldots ,i_p)\) 和 \(t = (j_1,\ldots ,j_q)\) ，则 \(s.t\) 是序列 \((i_1,\ldots ,i_p,j_1,\ldots ,j_q)\) ，共有 \(p + q\) 项。于是 \(e_{s,t} = e_s\otimes e_t\) ，从而由 (66) 可得

\[
\left\{ \begin{array}{l} e _ {s} \llcorner e _ {t} ^ {*} = 0 \quad \text { if   } s \text {   is   not   of   the   form   } t. u \\ e _ {t. u} \llcorner e _ {t} ^ {*} = e _ {u}. \end{array} \right.\tag{72}
\]

类似地，对称代数 \(\mathbf{S}(\mathbf{M})\) 以单项式 \(e^{\alpha}\) （ \(\alpha \in \mathbf{N}^{n}\) ）的集合为基（§ 6，第 6 号，定理 1），而 \(\mathbf{S}(\mathbf{M}^{*})\) 以单项式 \(e^{*\alpha}\) （ \(\alpha \in \mathbf{N}^{n}\) ）的集合为基；回顾（第 5 号）， \(u_{\alpha}\) （对 \(|\alpha| = k\) ）表示 \((\mathbf{S}^{k}(\mathbf{M}))^{*}\) 的基中与 \((e^{\alpha})_{|\alpha| = k}\) （ \(\mathbf{S}^{k}(\mathbf{M})\) 的基）对偶的元素；因此 \(u_{\alpha}\) （对 \(\alpha \in \mathbf{N}^{n}\) ）构成 \(\mathbf{S}(\mathbf{M})^{*\mathrm{gr}}\) 的一组基。由 \(e^{\beta}\) 所作的右内乘（在 \(\mathbf{S}(\mathbf{M})^{*\mathrm{gr}}\) 中）定义为乘以 \(e^{\beta}\) （在 \(\mathbf{S}(\mathbf{M})\) 中）的映射的转置，由此可知

\[
\left\{ \begin{array}{l l} u _ {\alpha} \llcorner e ^ {\beta} = 0 & \text { if } \quad \alpha \not \geqslant \beta \\ u _ {\alpha} \llcorner e ^ {\beta} = u _ {\alpha - \beta} & \text { if } \quad \alpha \geqslant \beta . \end{array} \right.\tag{73}
\]

类似地，由于 \(\mathbf{S}(\mathbf{M})\) 在此典范地与 \(\mathbf{S}(\mathbf{M})^{*\varepsilon r}\) 的分次对偶等同， \(i(u_{\beta})\) 是乘以 \(u^{\beta}\) （在 \(\mathbf{S}(\mathbf{M})^{*\varepsilon r}\) 中）的分次转置，因此由基 \((u_{\alpha})\) 的乘法表 (33)（第 5 号）推出

\[
\left\{ \begin{array}{l l} e ^ {\alpha} \llcorner u _ {\beta} = 0 & \text { if } \quad \alpha \not \geqslant \beta \\ e ^ {\alpha} \llcorner u _ {\beta} = (\beta , \alpha - \beta) e ^ {\alpha - \beta} & \text { if } \quad \alpha \geqslant \beta . \end{array} \right.\tag{74}
\]

至于 \(\mathsf{S}(\mathbf{M})\) 中元素被 \(\mathsf{S}(\mathbf{M}^*)\) 中元素所作的右内乘，此乘积的定义（第 9 号）以及第 5 号的公式 (34) 使我们能从 (74) 推出以下公式

\[
\left\{ \begin{array}{l l} e ^ {\alpha} \llcorner e ^ {* \beta} = 0 & \text { if } \quad \alpha \not \geqslant \beta \\ e ^ {\alpha} \llcorner e ^ {* \beta} = \frac {\alpha !}{(\alpha - \beta) !} e ^ {\alpha - \beta} & \text { if } \quad \alpha \geqslant \beta . \end{array} \right.\tag{75}
\]

对 \(\mathsf{S}(\mathbf{M}^{*})\) 中元素被 \(\mathsf{S}(\mathbf{M})\) 中元素所作的内乘，存在类似的公式：交换 M 与 \(M^{*}\) 的角色（因为此处 \(M^{**}\) 与 M 等同）。

注记。给定基 \((e_i)_{1\leq i\leq n}\) 使我们能够把代数 \(\mathsf{S}(\mathbf{M})\) 与多项式代数 \(\mathbf{A}[\mathbf{X}_1,\ldots ,\mathbf{X}_n]\) 等同（§ 6，第 6 号）；公式 (75) 表明，由 \(e^{*\alpha}\) 所作的内乘正是微分算子 \(\mathbf{D}^{\alpha} = \mathbf{D}_{1}^{\alpha_{1}}\mathbf{D}_{2}^{\alpha_{2}}\ldots \mathbf{D}_{n}^{\alpha_{n}}\) ，其中 \(\mathbf{D}_i = \partial /\partial \mathbf{X}_i\) （对 \(1\leqslant i\leqslant n\) ）（§ 10，第 11 号，例子）。

最后考虑外代数 \(\bigwedge (\mathbf{M})\) ，它以元素 \(e_{\mathbf{J}}\) 的集合为基，其中 \(\mathbf{J}\) 遍历区间 \([1, n]\) （属于 \(\mathbf{N}\) ）的子集的集合（§ 7，第 8 号，定理 1）；类似地， \(\bigwedge (\mathbf{M}^*)\) 以元素 \(e_{\mathbf{J}}^{*}\) 为基。由第 9 号的公式 (68) 可知

\[
\left\{ \begin{array}{l l} e _ {\mathrm{K}} ^ {*} \lrcorner e _ {\mathrm{J}} = 0 & \text { if } \quad \mathrm{K} \nsubseteq \mathrm{J} \\ e _ {\mathrm{K}} ^ {*} \lrcorner e _ {\mathrm{J}} = (- 1) ^ {p (p - 1) / 2} \rho_ {\mathrm{K}, \mathrm{J} - \mathrm{K}} e _ {\mathrm{J} - \mathrm{K}} & \text { if } \quad \mathrm{K} \subset \mathrm{J} \quad \text { and } \quad p = \operatorname{Card} (\mathrm{K}), \end{array} \right.\tag{76}
\]

（其中 \(\rho_{\mathbf{K},J - \mathbf{K}}\) 是由 § 7 第 8 号公式 (19) 定义的数。把 \(\mathbf{M}\) 和 \(\mathbf{M}^*\) 的角色互换，便有类似的公式。）

